\section{Machine Learning Fundamentals}\label{sec:bg:ml}

\Gls{ml} replaces hand-coded decision rules with parametric models fitted to data.
A model is a function \(\Enc\colon \setX \to \setY\) from an input space \(\setX\) to a task-specific output space \(\setY\).
Learning searches over parameters \(\bm{\theta}\) so that \(\Enc\) matches the training examples under a chosen loss.

The usual training objective is empirical risk minimization.
Given a dataset \(\setD = \{(\vect{x}_i, y_i)\}_{i=1}^{N}\) and a pointwise loss \(\ell\), the learner solves
\begin{equation}\label{eq:bg:erm}
  \hat{\bm{\theta}}
  =
  \argmin_{\bm{\theta}}\;
  \frac{1}{N}\sum_{i=1}^{N} \ell\bigl(\Enc(\vect{x}_i),\, y_i\bigr) + \Omega(\bm{\theta}),
\end{equation}
where \(\Omega\) is an optional regularizer.
Generalization is measured on held-out data.

In supervised learning each \(\vect{x}_i\) has a label \(y_i\), and \(\ell\) penalizes disagreement with that label.
In unsupervised learning there are no labels, and the aim is to expose structure, for example by clustering.
This thesis is supervised at training time: recordings carry per-pulse emitter and mode labels (\cref{sec:method:formulation,sec:method:data}).
At inference the labels are recovered by clustering embeddings, so the clustering tools of \cref{sec:bg:clustering} and the contrastive losses of \cref{sec:bg:dl} still matter.
\Cref{sec:bg:transformer} describes the encoder.
